% GENERATED FILE. Edit canonical lessons, then run npm run generate:books.
% canonical-source-hash: fnv1a64:a4dcb8add07f95aa
% canonical-lessons: PA-C55-patni, PA-C55-bachcha, PA-C55-munda, PA-C55-kuri, PA-C55-admi

\chapter{Wife, Child, Boy, Girl, Man}
\label{ch:pa-wife-child-boy-girl-man}
\hlchaptermodality{55}

\begin{chapteropening}
\textbf{By the end of this chapter:} \emph{I can say a wife, a child, a boy, a girl and a man.}

Complete the last lesson of chapter 55: I can say a wife, a child, a boy, a girl and a man.
\end{chapteropening}

\section[\texorpdfstring{\pa{ਪਤਨੀ} (patnī)}{ਪਤਨੀ (patnī)}]{\pa{ਪਤਨੀ} (patnī) — a wife}

\label{lesson:PA-C55-patni}



\begin{quote}
Before the new one: say the Punjabi for a grandmother, then the Punjabi for a husband.
\end{quote}

\subsection*{You'll want to know: \pa{ਪਤਨੀ}}
\textbf{\pa{ਪਤਨੀ}} — \emph{patnī} — \textquotedblleft{}a wife\textquotedblright{}.

\begin{cousinweb}
\textbf{\pa{ਪਤਨੀ}} (\emph{patnī}) is Sanskrit \emph{patnī}, \textquotedblleft{}mistress, wife\textquotedblright{}, taken over whole as a learned loan, a tatsama (CDIAL 7742).
\end{cousinweb}

\begin{grammarlens}[title={What you've built}]
One more word for this chapter.
\end{grammarlens}

\subsection*{Your turn}
\begin{itemize}
\raggedright
  \item \emph{Say it:} \emph{patnī}
  \item \emph{Say it:} \emph{patnī}, once more
  \item \emph{Say it:} say \emph{patnī}
  \item \emph{Recall:} say the Punjabi for a grandmother, then the Punjabi for a husband, then say \emph{patnī} again
\end{itemize}

\begin{tcolorbox}[breakable,colback=teal!4,colframe=teal!35!black,title={Before you move on}]
What does \pa{ਪਤਨੀ} mean? (\textquotedblleft{}A wife\textquotedblright{}.) Say it once more.
\end{tcolorbox}
\section[\texorpdfstring{\pa{ਬੱਚਾ} (bachchā)}{ਬੱਚਾ (bachchā)}]{\pa{ਬੱਚਾ} (bachchā) — a child}

\label{lesson:PA-C55-bachcha}



\begin{quote}
Before the new one: say the Punjabi for a husband, then the Punjabi for a wife.
\end{quote}

\subsection*{You'll want to know: \pa{ਬੱਚਾ}}
\textbf{\pa{ਬੱਚਾ}} — \emph{bachchā} — \textquotedblleft{}a child\textquotedblright{}.

\begin{cousinweb}
CDIAL lists \textbf{\pa{ਬੱਚਾ}} (\emph{bachchā}) under Sanskrit \emph{apatya}, \textquotedblleft{}offspring\textquotedblright{}, through Prakrit \emph{avacca} (CDIAL 424), and its addenda set the word beside Persian \emph{bachcha}, \textquotedblleft{}young (of humans or animals)\textquotedblright{} (CDIAL 14226). Many take it from Persian instead, so its origin is not settled.
\end{cousinweb}

\begin{grammarlens}[title={What you've built}]
One more word for this chapter.
\end{grammarlens}

\subsection*{Your turn}
\begin{itemize}
\raggedright
  \item \emph{Say it:} \emph{bachchā}
  \item \emph{Say it:} \emph{bachchā}, once more
  \item \emph{Say it:} say \emph{bachchā}
  \item \emph{Recall:} say the Punjabi for a husband, then the Punjabi for a wife, then say \emph{bachchā} again
\end{itemize}

\begin{tcolorbox}[breakable,colback=teal!4,colframe=teal!35!black,title={Before you move on}]
What does \pa{ਬੱਚਾ} mean? (\textquotedblleft{}A child\textquotedblright{}.) Say it once more.
\end{tcolorbox}
\section[\texorpdfstring{\pa{ਮੁੰਡਾ} (muṇḍā)}{ਮੁੰਡਾ (muṇḍā)}]{\pa{ਮੁੰਡਾ} (muṇḍā) — a boy}

\label{lesson:PA-C55-munda}



\begin{quote}
Before the new one: say the Punjabi for a wife, then the Punjabi for a child.
\end{quote}

\subsection*{You'll want to know: \pa{ਮੁੰਡਾ}}
\textbf{\pa{ਮੁੰਡਾ}} — \emph{muṇḍā} — \textquotedblleft{}a boy\textquotedblright{}.

\begin{cousinweb}
CDIAL lists \textbf{\pa{ਮੁੰਡਾ}} (\emph{muṇḍā}), \textquotedblleft{}boy\textquotedblright{}, under Sanskrit \emph{muṇḍa}, \textquotedblleft{}shaved, bald\textquotedblright{}, which it takes, probably, from Dravidian; it calls a Munda source less likely (CDIAL 10191).
\end{cousinweb}

\begin{grammarlens}[title={What you've built}]
One more word for this chapter.
\end{grammarlens}

\subsection*{Your turn}
\begin{itemize}
\raggedright
  \item \emph{Say it:} \emph{muṇḍā}
  \item \emph{Say it:} \emph{muṇḍā}, once more
  \item \emph{Say it:} say \emph{muṇḍā}
  \item \emph{Recall:} say the Punjabi for a wife, then the Punjabi for a child, then say \emph{muṇḍā} again
\end{itemize}

\begin{tcolorbox}[breakable,colback=teal!4,colframe=teal!35!black,title={Before you move on}]
What does \pa{ਮੁੰਡਾ} mean? (\textquotedblleft{}A boy\textquotedblright{}.) Say it once more.
\end{tcolorbox}
\section[\texorpdfstring{\pa{ਕੁੜੀ} (kuṛī)}{ਕੁੜੀ (kuṛī)}]{\pa{ਕੁੜੀ} (kuṛī) — a girl}

\label{lesson:PA-C55-kuri}



\begin{quote}
Before the new one: say the Punjabi for a child, then the Punjabi for a boy.
\end{quote}

\subsection*{You'll want to know: \pa{ਕੁੜੀ}}
\textbf{\pa{ਕੁੜੀ}} — \emph{kuṛī} — \textquotedblleft{}a girl\textquotedblright{}.

\begin{cousinweb}
CDIAL lists \textbf{\pa{ਕੁੜੀ}} (\emph{kuṛī}), \textquotedblleft{}girl, daughter\textquotedblright{}, under a reconstructed \emph{*kuḍa}, \textquotedblleft{}boy, son\textquotedblright{}, with \emph{*kuḍī}, \textquotedblleft{}girl, daughter\textquotedblright{}. It takes the word, probably, from Munda, comparing Santali \emph{kuṛi}, \textquotedblleft{}girl\textquotedblright{}; or else from Dravidian (CDIAL 3245).
\end{cousinweb}

\begin{grammarlens}[title={What you've built}]
One more word for this chapter.
\end{grammarlens}

\subsection*{Your turn}
\begin{itemize}
\raggedright
  \item \emph{Say it:} \emph{kuṛī}
  \item \emph{Say it:} \emph{kuṛī}, once more
  \item \emph{Say it:} say \emph{kuṛī}
  \item \emph{Recall:} say the Punjabi for a child, then the Punjabi for a boy, then say \emph{kuṛī} again
\end{itemize}

\begin{tcolorbox}[breakable,colback=teal!4,colframe=teal!35!black,title={Before you move on}]
What does \pa{ਕੁੜੀ} mean? (\textquotedblleft{}A girl\textquotedblright{}.) Say it once more.
\end{tcolorbox}
\section[\texorpdfstring{\pa{ਆਦਮੀ} (ādmī)}{ਆਦਮੀ (ādmī)}]{\pa{ਆਦਮੀ} (ādmī) — a man}

\label{lesson:PA-C55-admi}



\begin{quote}
Before the new one: say the Punjabi for a boy, then the Punjabi for a girl.
\end{quote}

\subsection*{You'll want to know: \pa{ਆਦਮੀ}}
\textbf{\pa{ਆਦਮੀ}} — \emph{ādmī} — \textquotedblleft{}a man\textquotedblright{}.

\begin{cousinweb}
\textbf{\pa{ਆਦਮੀ}} (\emph{ādmī}) is a loan from Arabic \emph{ādamī}, \textquotedblleft{}human, of Adam\textquotedblright{}, which reached Punjabi through Persian. Its Arabic root is \emph{ʾ-d-m} (Lane). It is \emph{Ādam}, \textquotedblleft{}Adam\textquotedblright{}, with the Arabic ending \emph{-ī}, \textquotedblleft{}belonging to\textquotedblright{}.
\end{cousinweb}

\begin{grammarlens}[title={What you've built}]
One more word for this chapter.
\end{grammarlens}

\subsection*{Your turn}
\begin{itemize}
\raggedright
  \item \emph{Say it:} \emph{ādmī}
  \item \emph{Say it:} \emph{ādmī}, once more
  \item \emph{Say it:} say \emph{ādmī}
  \item \emph{Recall:} say the Punjabi for a boy, then the Punjabi for a girl, then say \emph{ādmī} again
\end{itemize}

\begin{tcolorbox}[breakable,colback=teal!4,colframe=teal!35!black,title={Before you move on}]
What does \pa{ਆਦਮੀ} mean? (\textquotedblleft{}A man\textquotedblright{}.) Say it once more.
\end{tcolorbox}
